\documentclass[11pt]{article}
\usepackage[T1]{fontenc}
\usepackage[utf8]{inputenc}
\usepackage[a4paper,margin=1in]{geometry}
\usepackage{amsmath,amssymb}
\usepackage[draft]{graphicx}
\usepackage{hyperref}

\newcommand{\onecolumngrid}{}
\newcommand{\twocolumngrid}{}

\begin{document}

\section{\large{End Matter}}

\twocolumngrid

\emph{Appendix A: Berry phase Hamiltonian}.---In this section, we show how the combination of the Berry phase with $N_J$ fluctuations gives rise to the effective Hamiltonian in the main text for adiabatic $\delta$. 
In the limit of large $N$, the steady state in the spin-polarized phase becomes a pure state that is approximately an eigenstate of $\hat J^-$ with eigenvalue $\alpha = j o e^{i \phi}$, where $j \equiv N_J/2$ and $o \equiv |\Omega/\Omega_c| = |\Omega|/(j \sqrt{\Gamma^2 + 4 \chi^2})$, where we have included collective elastic interactions $\chi \hat J^+ \hat J^-$ in the Hamiltonian \cite{Somech2024}. The eigenvalue phase $\phi$ is determined by the phase of $\Omega$ and the value of $\chi/\Gamma$ but is otherwise unimportant in determining the Berry phase of the steady state, although $\chi$ will modify the Liouvillian spectrum and its eigenstates. These eigenstates become exact exponentially fast with increasing $N_J$ provided $\langle \hat J_z \rangle = \cos \theta_J \gtrsim N_J^{-1/3}$, i.e., the system is not too close to the critical point where finite-size effects play a strong role. 
 
The corresponding eigenstate can be expressed \cite{Somech2024}
\begin{subequations}
\begin{equation}
    |\alpha \rangle = \sum_m a_m |j,m \rangle, \quad
    a_m = a_{-j}e^{f_m} e^{-i (N_J/2 + m) \phi}, 
\end{equation}
\begin{equation}
\begin{aligned}
    f_m ~&= (m+N_J/2) \log (o N_J/2) ~- \\ & \sum_{k = -N_J/2}^{m-1} \frac{1}{2} \log \left[\left(\frac{N_J}{2}+1\right) \frac{N_J}{2}-(k+1)k\right],
\end{aligned}
\end{equation}
\end{subequations}
where $|N_J/2,m\rangle$ are collective Dicke states and $a_{-j}$ is determined by normalizing the state. We select the gauge by taking $a_{-j}$ to be independent of $\phi$, corresponding to $\phi$ information being contained in $|e\rangle$, consistent with the rotating frame used for $\hat H_\delta$. Variations in the phase of $\Omega$ correspond to variations in $\phi$. Thus, the Berry connection in the $\phi$ direction is 
\begin{equation}
\begin{aligned}
    \mathcal{A}_\phi &= \frac{1}{r \sin \theta_J} i \langle \alpha| \nabla_\phi|\alpha \rangle = \frac{1}{r \sin \theta_J} \sum_m |a_m|^2 \left( \frac{N_J}{2} + m \right) \\
    &= \frac{\langle \hat J_z \rangle}{r \sin \theta_J} = \frac{N_J(1- \cos \theta_J)}{2 r \sin \theta_J} + \mathcal{O}(1),
\end{aligned}
\end{equation}
where we parametrize the variation according to the steady state Bloch vector, and the Berry connection is the same as for $N_J$ non-interacting atoms pointing in the direction of the steady state up to higher-order corrections in $1/N_J$. This applies to $\mathcal{A}_{\theta_J}, \mathcal{A}_r$ as well, so only $\mathcal{A}_\phi$ is relevant for the Berry phase in the gauge we have chosen. As a result, the phase accumulation rate of our protocol when $\delta$ is applied is 
\begin{equation}
\begin{aligned}
    \partial_t \phi_B[N_J] &= \delta \frac{N_J (1-\cos \theta_J)}{2}  + \mathcal{O}(1)\\
    &= \delta \frac{N_J (1-\sqrt{1 - o^2})}{2}+ \mathcal{O}(1),
\end{aligned}
\end{equation}
where we have utilized the fact that $\cos \theta_J = \sqrt{1-o^2} + \mathcal{O}(1/N_J)$ for adiabatic $\delta$.

As in the main text, we proceed to include fluctuations in $N_J$ to identify the effective dynamics. As such, we take
\begin{equation}
    N_J \to \frac{N}{2} - \hat S_z, \qquad o \to \frac{N/2}{N/2-\hat S_z} \tilde{o},
\end{equation}
where we have expressed $o$ in terms of the mean-field value $\tilde{o}$ at $\langle \hat N_J \rangle = N/2$. We expand in powers of $\hat S_z/N$, leading to
\begin{equation}
\begin{aligned}
    \partial_t \phi_B = \tilde{\omega}_B \hat S_z + \frac{\delta \sin^2 \tilde{\theta}_J}{2 N \cos^3 \tilde{\theta}_J} \hat S_z^2 + \cdots,
    \end{aligned}
\end{equation}
where $\tilde{o} \equiv \sin \tilde{\theta}_J$ and we have dropped an unimportant constant term. This describes the phase evolution of each $N_J$ term in the superposition, Eq.~(\ref{eq:steadypsi}), and hence the effective Hamiltonian $\hat H_\text{eff} \equiv -\partial_t \phi_B$.
Thus, we see that $\omega_B$ corresponds to the rotations in phase on the $\hat S$ Bloch sphere, which is removed in an appropriate rotating frame. We can include non-adiabatic corrections by incorporating the effect of $\delta$ on $\cos \theta_J$.

\emph{Appendix B: Collective dephasing}.---In this section, we identify the strength of collective dephasing in the adiabatic $\delta$ limit in a similar fashion. This is accomplished by considering how the emission of light is modified in the presence of a detuning. To understand this, we recast our dynamics as
\begin{equation}
    \hat H = -\delta \hat N_e, \qquad \hat l = \hat J^- + i \Omega/\Gamma,
\end{equation}
where the Rabi frequency has been absorbed into the jump operator. This further clarifies how in the resonant model ($\delta = 0$) the steady state is an eigenvector of $\hat J^-$: the corresponding eigenvector is a dark state of the modified jump operator. As discussed in the main text, this shifted jump operator is proportional to the steady-state cavity coherence and thus captures the distinguishability of different $N_J$ terms by the environment.

The role of $N_J$ fluctuations in the detuning Hamiltonian are already captured in the previous analysis of the Berry phase since $\langle \hat N_e \rangle = N_J(1- \cos \theta_J)/2$, so we focus on the remaining jump operator. Here, we utilize the mean-field values at linear order in $\delta$, particularly $\sin \theta_J = \frac{\Omega}{N_J \Gamma/2}$,
\begin{equation}
\begin{aligned}
    \langle \hat l \rangle &= N_J \sin^2 \theta_J\left(\frac{\delta}{2 \Omega} \sec \theta_J - i \frac{N_J \Gamma}{4 \Omega} \right) + i \Omega/\Gamma \\
    &= \frac{\delta}{\Gamma} \frac{o}{\sqrt{1-o^2}}\\
    & = \frac{\delta}{N \Gamma} \left(N \tan \tilde{\theta}_J + \frac{2 \sin \tilde{\theta}_J}{ \cos^3 \tilde{\theta}_J } \hat S_z + \cdots\right) ,
    \label{eq:flucdeph}
\end{aligned}
\end{equation}
where the last step has utilized the same approach to incorporating $N_J$ fluctuations as for the Berry phase. Thus, we see there that the steady state is no longer a dark state, leading to $\hat S_z$-dependent emission of the light. The linear $\hat S_z$ term defines an effective collective dephasing due to the environment's ability to slowly gain information about $N_J$, analogous to a quantum non-demolition measurement \cite{Schleier-Smith2010a, Hosten2016, Cox2016, Barberena2023b}. The constant term has no effect on the dynamics and can be dropped.

\emph{Appendix C: Benchmarking squeezing dynamics}.---In this section, we benchmark our analytic expressions for the OAT and collective dephasing rates with exact numerics. For the exact numerics, we utilize a quantum trajectory approach. While the whole Hilbert space is $\mathcal{O}(N^3)$, we utilize two simplifications. First, the strong symmetry allows us to express $\hat J$ in terms of independent operators for each $N_J$ sector, reducing the size of the operators we consider. Second, for the initial equal superposition of $|{\downarrow}\rangle$ and $|{\uparrow}\rangle$, the distribution of $N_J$ is a binomial distribution centered around $N_J = N/2$. Therefore, we restrict the range of $N_J$ in the initial state to $N/2 \pm dN$. Additionally, we utilize the jump operator $\hat l = \hat J^- + i \Omega/\Gamma$ to best reflect the actual photon loss, which also minimizes the frequency of jumps at the steady state.

To calculate the spin squeezing in the system in the middle of the protocol, we consider a generalization of the Wineland squeezing parameter. Taking $|c \rangle$ to be the single-particle state associated with the Bloch vector and identifying the remaining orthogonal states $|j\rangle, |s\rangle$ states, we capture the fluctuations perpendicular to the Bloch vector via the operators
\begin{subequations}
    \begin{equation}
         \sum_i \frac{|c_i \rangle \langle j_i| + |j_i \rangle \langle c_i|}{2}, \qquad
         \sum_i \frac{|c_i \rangle \langle j_i| - |j_i \rangle \langle c_i|}{2i},
    \end{equation}
    \begin{equation}
        \sum_i \frac{|c_i \rangle \langle s_i| + |s_i \rangle \langle c_i|}{2}, \qquad
        \sum_i \frac{|c_i \rangle \langle s_i| - |s_i \rangle \langle c_i|}{2i},
    \end{equation}
\end{subequations}
which are analogous to a pair of fluctuations with $x$ and $y$ components. With a proper choice of $|j\rangle,|s\rangle$, these can be understood as fluctuations on the $\hat J, \hat S$ Bloch spheres, respectively. We define an associated covariance matrix $\mathcal{C}$, from which the generalized (anti-)squeezing is given by eigenvalues of
\begin{equation}
    \frac{ N \mathcal{C}}{\langle \hat N_c/2\rangle^2},
\end{equation}
where $\hat N_c/2 \equiv \sum_i |c_i \rangle \langle c_i|/2$ defines the multilevel Bloch vector length. 

In Fig.~\ref{fig:benchmark}, we compare the exact dynamics of the squeezing to Eq.~(\ref{eq:AEHP}) for $N=1000$, $dN = 80$, $ \delta = 0.05 N \Gamma, \Omega = 0.465 N \Gamma$, corresponding to $\cos \tilde{\theta}_J \approx 0.5$, with an expected relative ratio of OAT to collective dephasing rates of 0.44, i.e., in the limit of significant dephasing. We consider 500 trajectories with time steps $dt \leq (N \Gamma)^{-1}/250$, which we reduce by a factor of 5 briefly after each time the drive is turned on or off, limiting the quantum jump probability at a given step to under 5\%. For the analytics, we consider a fixed, steady-state value $\cos \tilde{\theta}_J \approx 0.5$. 

We find excellent agreement, with the analytics ($-7.3$ dB) appearing to underestimate the exact results (minimum: $-7.8$ dB; steady-state: $-7.6$ dB). We attribute this disagreement to non-adiabatic and finite-size effects that result in slight $\hat J$ correlations even after the drive is turned off, illustrated by the small nonzero steady-state $\langle e \rangle \approx 6 \times 10^{-3}$ after the drive is off. If instead of considering $\mathcal{C}$, we instead calculate the squeezing in the $|{\uparrow}\rangle,|{\downarrow}\rangle$ basis directly, we find a steady-state squeezing of $-7.1$ dB, which differs from the analytic result by the same amount that is lost during the dynamics after the drive is turned off.

\begin{figure}[h]
    \centering
    \includegraphics[scale=0.5]{DissSqBenchmark.pdf}
    \caption{(a) Comparison of squeezing dynamics for exact quantum trajectories (blue) to analytics (red) for $\delta = 0.05 N \Gamma, \cos \tilde{\theta}_J \approx 0.5$.  For the dark blue line, the drive is turned off at $\Gamma t = 0.17$ (star). (b) Excited state fraction $\langle e \rangle$ dynamics during the protocol.}
    \label{fig:benchmark}
\end{figure}

\end{document}
